\section{Justificativa}

A análise comparativa de diferentes estratégias de fatiamento permite compreender como suas características influenciam a utilização dos recursos disponíveis em diferentes condições da rede. \citeonline{gomes2020strategies} demonstram que estratégias distintas podem favorecer diferentes aspectos da infraestrutura, como disponibilidade de banda, saturação dos recursos e quantidade de clientes atendidos. Essa análise torna-se especialmente relevante em ambientes de demanda elástica, nos quais as necessidades de recursos podem variar ao longo do tempo. Nesse cenário, conhecer as vantagens e limitações de cada estratégia pode contribuir para a escolha de abordagens mais adequadas às diferentes condições de utilização da infraestrutura.

Em um ambiente de demanda elástica, a avaliação dessas estratégias sob diferentes condições de utilização torna-se importante, pois o comportamento da rede pode se modificar à medida que a demanda aumenta ou diminui. Uma estratégia que apresenta resultados satisfatórios em determinado nível de utilização pode não manter o mesmo comportamento diante de outras condições de demanda. Dessa forma, a realização de experimentos que considerem essas variações possibilita uma análise mais ampla do desempenho dos algoritmos, evitando que sua avaliação seja limitada a um único cenário.

A comparação entre diferentes algoritmos de fatiamento pode contribuir para uma melhor compreensão das características de cada abordagem e dos cenários nos quais sua utilização se mostra mais adequada. A identificação desses comportamentos permite observar não apenas qual estratégia apresenta melhores resultados em determinada condição, mas também quais fatores podem influenciar seu desempenho. Dessa forma, estudos comparativos podem contribuir para o desenvolvimento e aprimoramento de mecanismos de alocação de recursos em ambientes de demanda elástica.